\section{Zadanie 3}
\paragraph{Model programowania liniowego}
\begin{enumerate}
\item Parametry \\
m - liczba różnych zasobów, i $\in$ N \\ 
n - liczba różnych produktów, j $\in$ N \\
$c_{i}^{max}$ - Maksymalna przepustowość dla zasobu i \\
$A_{ij}$ - Użycie zasobu i przez produkt j \\ 
$q_j$ - Limit produkcji dla produktu j, po którym zmniejsza się przychód na jednostce \\
$p_j$ - Normalny przychód na pojedyńczej sztuce produktu j \\ 
$p_j^{disc}$ - Pomniejszony przychód na sztuce produktu j 

\item Zmienne decyzyjne 
$x_j$ - Wolumin produkcji produktu j 
\item Zmienne pomocnicze 
$r_j(u)$ - Funkcja wyliczająca przychód dla produktu j bazując na poziomie produkcji u \\ 
$c_i$ - Całkowite zużycie zasobu typu i 

\item Funkcja celu 
\[ \max Z = \sum_{j=1}^n r_j u_j \] 
\[ r_j(u) = \{ \]
\[    p_{j} u \qquad \qquad \qquad \qquad u <= q_j  \]
\[    p_j q_j + p_{j}^{disc}(u-q_j), \qquad u => q_j \]
\[ \} \]

\item Ograniczenia 
\setcounter{equation}{0}
Całkowite zużycie zasobów nie może przekraczać dostępnośli dla żadnego zasobu
\begin{equation}
\sum_{j=1}^n A_{ij} \cdot x_j \leq c_i^{max}, \qquad i \in N
\end{equation}
Produkcja każdego produktu j, jest powiżana z właściwym $u_j$
\begin{equation}
u_j = x_j, \qquad j \in N
\end{equation}
Wolumin produkcji nie może być mniejszy od zera
\begin{equation}
x_j \geq 0, \qquad j \in N
\end{equation}


\end{enumerate}
